% !TeX root = ../../main.tex
% chapters/5_mobiwac/02_related.tex -- one section of Chapter 5 (MobiWac 2026, submitted, under review).
% Split out of chapters/5_mobiwac.tex on 2026-07-28 (author request: the paper
% chapters were single long files, while the originals under articles/[mobiwac]/src/sections/
% are one file per section). The split is PURELY MECHANICAL: content was cut at
% \section boundaries and not edited. The chapter file now \inputs these in order.
% Editing rules are unchanged and still governed by the chapter's status above:
% see NORTH_STAR section 4 before changing any sentence here.
\section{Background and Related Work}
\label{sec:mobiwac:related}

\subsection{The MTLnet framework and the representation diagnosis}
\label{sec:mobiwac:related-recap}
% [NEEDS SIGN-OFF: raised v1 assembly, 2026-07-23 | new-to-chapter recap subsection (the bridging device mandated by
% NORTH_STAR section 3: Ch.5 recaps BOTH the Ch.3 artifact and the Ch.4 finding). Content
% from the approved spine (NORTH_STAR sections 2 and 6) and the GLOSSARY model-lineage table;
% no numbers are quoted.]
Two chapters of this dissertation precede this study, and this chapter answers the
question that they leave open. Chapter~\ref{ch:cbic} introduced MTLnet, the first
joint model for this task pair: task-specific encoders feed shared layers
conditioned on task identity, and a gradient-balancing method weights the two tasks
during training~\cite{silva2025mtlnet}. On a place-level embedding with hard
parameter sharing, that joint model did not outperform the dedicated single-task
models, the conclusion of the time for that configuration.
Chapter~\ref{ch:courb} kept the MTLnet architecture unchanged and replaced its
single 64-dimensional place-level input with decomposed spatial, temporal, and
categorical encoders (ST-MTLNet); the category score rose at each of the three
states studied, so the input representation, not the sharing architecture, was
identified as the bottleneck~\cite{paiva2026stmtlnet}. This chapter carries that diagnosis one step
further: the representation moves from the place level to the check-in level, the
sharing topology is redesigned around it, and the earlier negative conclusion is
revisited under this new configuration.

\subsection{From place embeddings to per-visit embeddings}
\label{sec:mobiwac:related-embeddings}
A place embedding transforms a POI into a vector so that similar or nearby
places sit close together in the vector space.
Deep Graph Infomax (DGI)~\cite{velickovic2019deep}, a general self-supervised
method for graphs, learns such vectors on a place network linked by
similarity, time, and distance, training them to tell the real network from a
copy with shuffled node features.
Hierarchical Graph Infomax (HGI) \cite{huang2023hgi} adds geographic structure,
organizing the network as place, then region, then city. Both produce one vector per place, so two visits to the same coffee shop look
identical to the model.
Contextual check-in representations instead give each visit its own vector: CTLE
\cite{lin2021ctle}, the closest example, learns a vector per visit by masking and
reconstructing parts of a user's check-in sequence. Our
representation, Check2HGI, also works at the check-in level but through a
different construction. CTLE is a sequence model, a Transformer that reads the
check-in sequence itself; Check2HGI remains a network model, keeping the
hierarchical place-to-region-to-city network and the same infomax objective,
now extended one level deeper, from individual places down to individual
check-ins (Figure~\ref{fig:mobiwac:dataflow}). The order of a user's visits still enters
the network, as links between consecutive check-ins
(Section~\ref{sec:mobiwac:method-rep}).

The novelty is this specific combination: per-visit context inside a
hierarchical graph-infomax representation. We compare against CTLE and a feature-concatenation control to separate the
combination from contextualization alone and from feature injection.

\subsection{Predicting the next category and the next region}
\label{sec:mobiwac:related-tasks}
A large body of work predicts the next place that a user will visit, the exact POI,
with recurrent and attention models such as DeepMove \cite{feng2018deepmove}, STAN
\cite{luo2021stan}, and GETNext \cite{yang2022getnext}. In contrast, we target two properties of the next visit, its category and its
region, rather than the exact place. We choose them because they are
what a mobility-aware service can act on, not to make the task easier; the region
alone spans hundreds to several thousand classes. Predicting
over a partition of the map is also the standard formulation in the human-mobility
literature, with a grid cell as the target~\cite{luca2021mobilitysurvey}; our
next-region task substitutes official neighborhood-scale units for grid cells. Our earlier
work~\cite{silva2025mtlnet} paired static category classification with next-category prediction
and found no consistent multitask gain; this chapter introduces the next-region task and adds the check-in-level
representation, on which sharing helps instead of hurting (Section~\ref{sec:mobiwac:results-part2}).

The field
increasingly models several granularities at once; in those systems, category and
region are auxiliary signals that help a primary next-place task
(MCMG \cite{sun2024mcmg}, HMT-GRN \cite{Lim2022}).
We study the pair as the object itself. To our knowledge, fine-grained region as an
end target of equal standing, rather than an auxiliary coarse grid cell, is
underexplored.
The nearest exceptions do not study our exact pairing:
DRRGNN~\cite{zhu2022drrgnn} forecasts a person's next activity region jointly
with a mobility-intention label, but over regions discovered per person rather
than a fixed citywide partition; a generative recommender~\cite{sun2025kgtb}
predicts category and region together, but only as auxiliary steps toward its
next-place ranking.

\subsection{Multitask learning for mobility}
\label{sec:mobiwac:related-mtl}
MCARNN \cite{Liao2018} jointly predicts activity and
location, and a recent hierarchy-aware model continues the
pattern \cite{wang2025hamtl}; in both, the location target is
the exact place. The closest alternative to our setting is
the cascade: predict the category first, then the place given the
category \cite{ye2013nextmove}. CSLSL
\cite{huang2024cslsl} is its current form, predicting in a chain (when, then what,
then where), with location as the primary output and category as an
intermediate step; CatDM \cite{yu2020catdm} likewise uses the predicted
category only to filter candidate places. We instead predict
category and region in parallel, as end targets of equal standing in one model
with a single forward pass, and we drop the next-place target entirely, so
neither task is an intermediate step toward a third. Since the cascade is the pattern that the field uses, we test the
choice directly (Section~\ref{sec:mobiwac:results-part2}). On optimization, we are conservative by design: joint
training with a fixed loss weighting is standard practice
\cite{caruana1997multitask}, not itself our contribution.
Reports find that gradient-balancing methods, which re-weight
the tasks' updates during training (PCGrad \cite{yu2020pcgrad}, Nash-MTL \cite{nash}), rarely improve on a
well-tuned fixed weighting with two tasks \cite{xin2022domtl}.
We confirm this at scale: of nineteen loss and gradient
balancers screened at their default configurations at a single seed on two datasets,
Alabama and Florida, including the two methods named above, none improved on a tuned fixed
task weighting across both tasks and both datasets. Two exceed equal weighting on next-category at Alabama, Nash-MTL by $0.68$
points and scale normalization by $0.19$, and neither holds that position elsewhere:
Nash-MTL loses it on the second dataset, and scale normalization gains on next-category
there while collapsing on next-region.
% [REV-011, author ruling 2026-07-24: "use the 'at their default configurations'"]
% Applied as written. Evidence check against docs/results/mtl_improvement/T4_audit_and_verdict.md:
%  - Nash-MTL (named): correctly wired and adapting, a VALID test (T4:37-39). Sound.
%  - PCGrad (named): one of three gradient-surgery methods that "never validly tested individually"
%    under the dual tower -- the private region tower sits outside shared_parameters(), so the
%    surgical gradient never reaches >80% of the region pathway and the method collapses to
%    approximately equal weighting (T4:26-31).
% RECOMMENDATION (agent, NOT applied -- the author's instruction governs): the new qualifier does
% not in fact cover PCGrad. Its exclusion is a WIRING result, invariant to configuration, so
% "at their default configurations" answers the DWA/FairGrad/GradNorm objection but not this one.
% The minimal repair that keeps the author's phrase is to drop the two words "PCGrad
% \cite{yu2020pcgrad}, " from the citing sentence above and let Nash-MTL carry the named evidence
% (the \cite{xin2022domtl} sentence already reports the literature claim in general terms).
% [round6, COD-005 residue] SCOPE ADDED, and this comment supersedes the "Also unstated" note that
% stood here. The sentence stated neither the repetition count nor which two datasets. Source of
% record, docs/results/mtl_improvement/T4_audit_and_verdict.md, opened this session: the
% evidence-strength banner at :8-10 reads "What actually ran: the full screen at registry DEFAULTS,
% seed 0, AL+FL", and the "Caveats / scope" section at :111-112 states it a second time,
% "Single-seed (seed0) screen + corrected re-run, FL+AL". The AL+FL expansion is not inferred from the
% abbreviation: T4_full_screen.json has exactly two top-level keys, alabama and florida, and each
% holds exactly the nineteen arms this sentence counts (aligned_mtl, bayesagg_mtl, cagrad, db_mtl, dwa,
% equal_weight, excess_mtl, fairgrad, famo, go4align, gradnorm, nash_mtl, pcgrad, scale_norm,
% scheduled_static, static_weight, stch, uncertainty_weighting, uw_so). Counted from the JSON, not
% from prose.
% "including the two named above" was ALSO ambiguous and is now "the two methods named above": the
% nearest antecedent is the pair of balancers PCGrad and Nash-MTL in the preceding sentence, not a pair
% of datasets. Measured 2026-07-28: neither state was named anywhere before this sentence in this file
% (the first "Alabama" was the next sentence, the first "Florida" was twenty-four lines below), so a
% reader could not have resolved "the two" to the datasets even if that had been intended. An earlier
% anchor table recorded this site as already naming the two datasets; that reading was wrong and is
% corrected here.
% The literal string "seed 0" stays out of the prose: it is a fact about the run configuration rather
% than a reported parameter, the same treatment the fold seed gets in Chapter 4, and
% articles/[mobiwac]/GLOSSARY.md:113 gives "a single seed" as the phrasing while banning the compound
% "single-seed states". No number, count or verb changed.
% PARITY: the same edit is applied in articles/[mobiwac]/src/sections/02_related.tex, since Chapter 5
% is under review and the two texts stay identical; logged in articles/[mobiwac]/ERRATA.md rather than
% in Appendix B.
The reason is visible in the gradients. Measured during development on the same
joint architecture (on an earlier preparation of the data), the cosine
similarity between the next-category and next-region updates on the shared
trunk averages $+0.001$ across training (four seeds each on four Gowalla
states: Alabama, Arizona and Florida, which are three of the five United States datasets reported
here, and Georgia, which this dissertation does not otherwise use; the largest per-dataset mean in
absolute value is $+0.0032$). A balancer therefore has no
conflict to resolve, consistent with the mechanism test in
Section~\ref{sec:mobiwac:results-part2}; this is a finding for this pair of tasks, not a
general rule. Appendix~\ref{apx:cosine} of this dissertation measures the same quantity on
the final model across seven datasets and reaches the same conclusion.
% [round9c, PENDENCIAS 2.9 + 2.23 (readability report R-6), 2026-07-30] ONE POINTER ADDED, and this is
% the DISSERTATION-ONLY sentence the two trees deliberately do not share. Errata row in Appendix B,
% Table B.5.
% WHY IT IS HERE. R-6 measured that the cosine finding reaches the reader at three separate places and
% that only one of them signposted the appendix: a reader who met the +0.001 figure in this chapter,
% saw it rests on four seeds and four states, and wanted the fuller measurement had no route to it,
% while a reader who arrived at the appendix having read two versions of the finding could take it for
% a third restatement. The report closed its own recommendation as unavailable ("Ch.5 is under review
% and its body is deliberately untouched, so the natural fix is closed") and asked the author instead.
% WHY IT IS AVAILABLE NOW. Item 2.9: he opened both Chapter 5 trees this round ("como estamos um fase
% de revisao no mobiwac, conseguimos mandar uma revisao ainda"). That permission is what makes the
% sentence possible, and item 2.9 asks whether the appendix's seven-dataset result warrants a mention
% here. It does, for the reason R-6 gives.
% WHY IT IS NOT MIRRORED INTO THE PAPER, which is the one exception to keeping the two texts
% identical. The sentence cites a dissertation appendix, so it cannot be folded into a paper under
% review at all: the referent does not exist there. That is the same class as the two corrections
% already in Table B.5, whose banner reads "each depends on material the article has no room for".
% Appendix B records it for exactly that reason, so the divergence is declared rather than silent.
% NO NEW CLAIM AND NO NUMBER. The appendix's conclusion is the same as this paragraph's, which is the
% point of pointing at it; the +0.001 figure and its four-state scope are untouched.
% [v2 review F-03, 2026-07-26] SCOPE CORRECTION (persona 10). The sentence said "three of our six
% datasets". The measurement pools FOUR Gowalla states -- alabama, arizona, georgia, florida --
% per docs/results/mtl_improvement/R0_matched_metric_bar.json (its "states" object has exactly those
% four keys, each with the four seeds 0/1/7/100) and scripts/.../plot_grad_cosine.py:19-24, whose
% STATE_STYLE map names the same four. Georgia is NOT one of the six datasets this dissertation
% reports, so "three of our six" both undercounted the pool and implied every state in it is one we
% report. Read from the JSON, not recomputed. The correction is in the author's favour: the finding
% replicates on a state the dissertation never reports.
% [round6, 2026-07-28] PARITY CLOSED. That F-03 correction had been applied on the dissertation side
% only, so until now the submitted paper still read "three of our six datasets" two lines from the
% sentence above, while this chapter named the four states. Under the Chapter 5 regime the two texts
% must be identical, so the dissertation's wording was carried into
% articles/[mobiwac]/src/sections/02_related.tex this round, with "this dissertation" rendered as
% "this study" there. WHICH TEXT WAS RIGHT: this one. I re-read the source myself rather than trusting
% the earlier correction. R0_matched_metric_bar.json has exactly four keys under "states" -- alabama,
% arizona, georgia, florida -- and "seeds": [0, 1, 7, 100]; scripts/mtl_improvement/plot_grad_cosine.py
% reads its per-state run directories out of that same JSON and its STATE_STYLE map names the same
% four. So four Gowalla states at four seeds each, of which three are among the five United States
% datasets this document reports. "Three of our six datasets" and "four Gowalla states" are different
% claims about the same measurement, and only the second one matches the artifact. Logged in
% articles/[mobiwac]/ERRATA.md.
% [round5, author decision 2.4] The text read as a two-optimizer test when the sweep was much
% wider. Source of record: docs/results/mtl_improvement/T4_full_screen.json, 19 methods x 2 states:
%   uncertainty_weighting, equal_weight, static_weight, scale_norm, scheduled_static, uw_so, dwa,
%   famo, gradnorm, db_mtl, cagrad, nash_mtl, aligned_mtl, pcgrad, fairgrad, stch, go4align,
%   bayesagg_mtl, excess_mtl.
% AL next-category macro-F1 vs equal_weight 53.57: TWO methods are above it, nash_mtl 54.25 (+0.68)
% and scale_norm 53.76 (+0.19). [correction] An earlier version of this line and of the sentence said
% nash_mtl was the ONLY one, while this very comment listed scale_norm as also above -- a claim that
% contradicted its own ledger. Caught in review. Full picture, all four cells:
%   nash_mtl:   AL cat 54.25 (+0.68) reg 62.77 (+0.24) | FL cat 71.61 (-0.19) reg 73.04 (-0.02)
%   scale_norm: AL cat 53.76 (+0.19) reg 61.96 (-0.57) | FL cat 72.12 (+0.33) reg 35.47 (-37.59)
% nash_mtl is the only one above on BOTH tasks at AL; scale_norm is the only one above on category
% at BOTH datasets, and its region score collapses. So no method wins everywhere, which is why
% the sentence says "across both tasks and both datasets" and then names both exceptions with
% what each one gives up. PCGrad 53.11 (-0.46) is kept by name at the
% author's request so an MTL reviewer sees it was tried.
% [NEEDS SIGN-OFF: raised round 5, 2026-07-27 | AUTHOR -- "nineteen", "+0.68" and "+0.19" are new numbers in the chapter.]

\begin{figure}[htbp]
\centering
% [round6, 2026-07-28] COD-017 in-figure text size. WRITING_LAW 5 requires in-figure text near
% body size. Measured with FPDFText_GetFontSize (not glyph bounding boxes): body is 11.96 pt and
% this diagram's labels were 6.97 pt, that is 58 percent of body, at 350 glyphs below 9 pt on
% rendered page 62. The figure is a vector PDF whose natural size is 391.170 by 177.718 pt, so it
% was being placed at natural size inside a 455.0 pt text block and left 63.8 pt of the block
% unused. Scaling the placement to the full text block scales the type with it. Regenerating the
% TikZ source at a larger base font was tried first and rejected: at 11pt the labels reach 7.97 pt
% natural but the node text then overruns its own box borders (measured on the sibling figure, the
% shared-trunk title protrudes 0.50 pt past the fill), which changes the figure rather than its
% size. Nothing inside the figure is altered here; only its printed scale.
\includegraphics[width=\textwidth]{figures/mobiwac/fig1_dataflow.pdf}
\caption{From a check-in sequence to two predictions: each visit is embedded in a four-level graph
(check-in, place, region, city) trained without task labels. The graph yields two sets of
vectors, one per visit and one per region; sliding nine-visit windows of each feed a single
model that outputs the next category and the next region.}
\label{fig:mobiwac:dataflow}
\end{figure}
